% Part: computability
% Chapter: recursive-functions
% Section: primitive-recursion

\documentclass[../../../include/open-logic-section]{subfiles}

\begin{document}

\olfileid[pt]{cmp}{rec}{pre}
\olsection{Recursão primitiva}

Uma característica dos números naturais é que todo número natural pode ser
alcançado a partir de~$0$ aplicando-se um número finito de vezes a operação
sucessor~$+1$---todo número natural é~$0$ ou o sucessor de \dots{} o
sucessor de~$0$. Uma maneira de especificar uma função
$h\colon\Nat \to \Nat$ que faça uso desse fato é a seguinte: (a)~especificar
o valor de $h$ para o argumento~$0$ e (b)~especificar também como, dado o
valor de $h(x)$, calcular o valor de $h(x+1)$. Com efeito, (a) nos diz diretamente
qual é $h(0)$, de modo que $h$ fica definida para~$0$. Usando então a
instrução dada por (b) para $x=0$, podemos calcular $h(1) = h(0+1)$ a partir
de~$h(0)$. Usando a mesma instrução para $x=1$, calculamos
$h(2) = h(1+1)$ a partir de~$h(1)$, e assim por diante. Para cada número
natural~$x$, acabaremos chegando à etapa em que definimos $h(x)$ a partir de
$h(x+1)$; portanto, $h(x)$ fica definida para todo $x \in \Nat$.

Por exemplo, suponha que especifiquemos $h\colon \Nat \to \Nat$ pelas duas
equações seguintes:
\begin{align*}
h(0) & =  1\\
h(x+1) & =  2 \cdot h(x)
\end{align*}
Se já sabemos multiplicar, essas equações nos dão as informações exigidas
por (a) e~(b) acima. Aplicando sucessivamente a segunda equação, obtemos
\begin{align*}
  h(1) & = 2\cdot h(0) = 2,\\
  h(2) & = 2\cdot h(1) = 2\cdot 2,\\
  h(3) & = 2 \cdot h(2) = 2\cdot 2 \cdot 2,\\
  & \vdots
\end{align*}
Vemos que a função~$h$ que especificamos é $h(x) = 2^x$.

A característica mencionada dos números naturais garante que há uma única
função~$h$ que satisfaz esses dois critérios. Um par de equações como esse é
chamado de \emph{definição por recursão primitiva} da função~$h$. O termo
``recursão'' é usado porque definimos~$h$ recursivamente, O.s.\@, a própria
$h$ ocorre no lado direito da definição, especificamente na segunda equação.
A recursão é ``primitiva'' porque, ao definir $h(x+1)$, usamos apenas o valor
$h(x)$, O.s.\@, o valor imediatamente anterior. Essa é a maneira mais simples
de definir recursivamente uma função em~$\Nat$.

Podemos definir por recursão primitiva até mesmo funções mais fundamentais,
como a adição e a multiplicação. Nesses casos, porém, as funções em questão
têm $2$ argumentos. Fixamos uma das posições de argumento e usamos a outra
para a recursão. Por exemplo, para definir $\Add(x, y)$, podemos fixar~$x$ e
definir primeiro o valor para $y=0$ e depois o valor para $y+1$ em termos
de~$y$. Como $x$ é fixo, ele aparecerá nos lados esquerdo e direito das
equações definidoras.
\begin{align*}
\Add(x,0) & =  x\\
\Add(x,y+1) & =  \Add(x,y)+1
\end{align*}
Essas equações especificam o valor de $\Add$ para todos os valores de $x$
\emph{e}~$y$. Para encontrar $\Add(2,3)$, por exemplo, aplicamos as equações
definidoras com $x = 2$: usamos a primeira para obter $\Add(2,0) = 2$ e, em
seguida, a segunda para obter sucessivamente $\Add(2,1) = 2 + 1 = 3$,
$\Add(2, 2) = 3 + 1 = 4$ e $\Add(2, 3) = 4 + 1 = 5$.

Na definição de $\Add$, usamos $+$ no lado direito da segunda equação, mas
apenas para somar~$1$. Em outras palavras, usamos a função sucessor
$\Succ(z) = z+1$ e a aplicamos ao valor anterior $\Add(x,y)$ para definir
$\Add(x,y+1)$. Assim, podemos pensar na definição recursiva como dada em
termos de uma única função que aplicamos ao valor anterior. Contudo, não há
problema---e às vezes é necessário---em permitir que a função dependa não
apenas do valor anterior, mas também de $x$ e~$y$. Considere:
\begin{align*}
  \Mult(x,0) & =  0 \\
  \Mult(x,y+1) & =  \Add(\Mult(x,y),x)
\end{align*}
Essa é uma definição por recursão primitiva de uma função $\Mult$: aplicamos
a função $\Add$ tanto ao valor anterior $\Mult(x,y)$ quanto ao primeiro
argumento~$x$. Ela também define a função~$\Mult(x,y)$ para todos os
argumentos $x$ e~$y$. Por exemplo, $\Mult(2,3)$ é determinado calculando-se
sucessivamente $\Mult(2,0)$, $\Mult(2,1)$, $\Mult(2,2)$ e~$\Mult(2,3)$:
\begin{align*}
  \Mult(2,0) & = 0\\
  \Mult(2,1) & = \Mult(2,0+1) =
  \Add(\Mult(2,0), 2) = \Add(0, 2) = 2\\
  \Mult(2,2) & = \Mult(2,1+1) =
  \Add(\Mult(2,1), 2) = \Add(2, 2) = 4\\
  \Mult(2,3) & = \Mult(2,2+1) =
  \Add(\Mult(2,2), 2) = \Add(4, 2) = 6
\end{align*}

O padrão geral é, então, o seguinte: para dar uma definição por recursão
primitiva de uma função $h(x_0, \dots, x_{k-1}, y)$, fornecemos duas equações.
A primeira define o valor de $h(x_0, \dots, x_{k-1}, 0)$ sem fazer referência
a~$h$. A segunda define o valor de $h(x_0, \dots, x_{k-1}, y+1)$ em termos de
$h(x_0, \dots, x_{k-1}, y)$, dos outros argumentos $x_0$,
\dots,~$x_{k-1}$ e de~$y$. Nessa segunda equação, só se pode usar o valor
imediatamente anterior de~$h$. Se pensarmos nas operações dadas pelos lados
direitos dessas duas equações como funções $f$ e~$g$, o padrão geral para
definir uma nova função~$h$ por recursão primitiva é este:
\begin{align*}
  h(x_0, \dots, x_{k-1}, 0) & = f(x_0, \dots, x_{k-1})\\
  h(x_0, \dots, x_{k-1}, y+1) & =
  g(x_0, \dots, x_{k-1}, y,\\
  &\qquad h(x_0, \dots, x_{k-1}, y))
\end{align*}  
No caso de $\Add$, temos $k=1$, $f(x_0) = x_0$ (a função identidade) e
$g(x_0, y, z) = z + 1$ (a função de $3$ argumentos que devolve o sucessor
de seu terceiro argumento):
\begin{align*}
  \Add(x_0, 0) & = f(x_0) = x_0\\
  \Add(x_0, y+1) & = g(x_0, y, \Add(x_0, y)) =
  \Succ(\Add(x_0, y))
\end{align*}
No caso de $\Mult$, temos $f(x_0) = 0$ (a função constante que sempre devolve
o valor~$0$) e $g(x_0, y, z) = \Add(z,x_0)$ (a função de $3$ argumentos que
devolve a soma de seu último e de seu primeiro argumento):
\begin{align*}
  \Mult(x_0,0) & =  f(x_0) = 0 \\
  \Mult(x_0,y+1) & =  g(x_0, y, \Mult(x_0,y)) =
  \Add(\Mult(x_0,y), x_0)
\end{align*}

\end{document}
